\section{Spectroscopic Constraints}
    \label{sec:spectroscopy}
    
    \subsection{Motivation}
        
        Any extension of atomic structure must remain consistent with high-precision spectroscopic measurements.
        
        \textbf{[ORTHODOX]}:
        Atomic energy levels are experimentally verified to extremely high precision.
        
        \textbf{Requirement:}
        \begin{align}
            |\delta E_n| \ll |E_n|
            \label{eq:spectro_constraint}
        \end{align}
        
        where $\delta E_n$ represents any SST-induced correction.
    
    \subsection{Internal Mode Hypothesis}
        
        We consider the possibility that atomic bound states possess an additional internal excitation sector associated with vortex filament dynamics.
        
        \textbf{[SPECULATIVE]}:
        Each bound state carries an internal mode spectrum $\{E_{m,n}\}$.
        
        The effective energy becomes:
        
        \begin{align}
            E_n^{\text{eff}} = E_n^{(0)} + \delta E_n
            \label{eq:effective_energy}
        \end{align}
    
    \subsection{Generic Coupling Structure}
        
        We parameterize the coupling between internal modes and observable energy levels:
        
        \begin{align}
            |\delta E_n| \leq \lambda_n \sum_m E_{m,n} \, \nu_{m,n}
            \label{eq:coupling_bound}
        \end{align}
        
        where:
        \begin{itemize}
            \item $\lambda_n$ is the coupling efficiency
            \item $\nu_{m,n}$ is the occupation factor
        \end{itemize}
        
        \textbf{[DERIVED]}: This form captures a general upper bound independent of microscopic details.
    
    \subsection{Ungapped Spectrum: Instability}
        
        If the internal spectrum is ungapped:
        
        \begin{itemize}
            \item arbitrarily low-energy excitations exist
            \item no universal suppression mechanism applies
        \end{itemize}
        
        \textbf{Consequence:}
        \begin{align}
            \delta E_n \not\to 0 \quad \text{generically}
        \end{align}
        
        \textbf{[CRITICAL]}:
        An ungapped internal sector is generically incompatible with spectroscopy unless:
        
        \begin{itemize}
            \item $\lambda_n \ll 1$ (extremely weak coupling), or
            \item the density of states is highly suppressed
        \end{itemize}
    
    \subsection{Gapped Spectrum: Decoupling Mechanism}
        
        We introduce an excitation gap $\Delta_K$:
        
        \begin{align}
            E_{m,n} \geq \Delta_K
        \end{align}
        
        Then occupation factors follow Boltzmann suppression:
        
        \begin{align}
            \nu_{m,n} \sim \exp\left(-\frac{\Delta_K}{k_B T_{\text{eff}}}\right)
        \end{align}
        
        Substituting into Eq.~\ref{eq:coupling_bound}:
        
        \begin{align}
            \delta E_n \propto \exp\left(-\frac{\Delta_K}{k_B T_{\text{eff}}}\right)
            \label{eq:boltzmann_suppression}
        \end{align}
        
        \textbf{[DERIVED]}:
        Exponential suppression ensures compatibility with spectroscopy.
    
    \subsection{Physical Origin of the Gap}
        
        \textbf{[SPECULATIVE]}:
        The gap $\Delta_K$ may arise from:
        
        \begin{itemize}
            \item geometric constraints on filament curvature
            \item torsional rigidity
            \item topological restrictions
            \item core structure of the vortex
        \end{itemize}
    
    \subsection{Order-of-Magnitude Estimate}
        
        Using filament parameters:
        
        \begin{align}
            \Delta_K \sim \frac{\Gamma^2}{4\pi \xi L}
        \end{align}
        
        where:
        \begin{itemize}
            \item $\xi$ is the core scale
            \item $L$ is the loop length
        \end{itemize}
        
        For typical SST scales:
        
        \begin{align}
            \Delta_K = \mathcal{O}(10^2 - 10^3 \, \text{eV})
        \end{align}
        
        \textbf{[DERIVED]}:
        This places the internal sector well above atomic energy scales ($\sim$ eV).
    
    \subsection{Consistency Condition}
        
        To preserve atomic physics:
        
        \begin{align}
            \frac{\Delta_K}{k_B T_{\text{eff}}} \gg 1
        \end{align}
        
        \textbf{Interpretation:}
        The internal sector is effectively frozen at atomic energies.
    
    \subsection{Observable Consequences}
        
        \textbf{[SPECULATIVE]}:
        Residual effects may include:
        
        \begin{itemize}
            \item tiny level shifts
            \item suppressed line broadening
            \item weak temperature dependence
        \end{itemize}
    
    \subsection{Falsifiability}
        
        The theory is falsifiable via:
        
        \begin{itemize}
            \item detection of unexplained spectral shifts
            \item absence of expected suppression
            \item deviation from Boltzmann scaling
        \end{itemize}
    
    \subsection{Integration with SST Framework}
        
        The spectroscopic constraint enforces:
        
        \begin{itemize}
            \item compatibility with the Atomic Bridge (Section~\ref{sec:atomic})
            \item constraints on delay-induced excitations (Section~\ref{sec:delay})
            \item bounds on allowable vortex configurations
        \end{itemize}
    
    \subsection{Canonical Statement}
        
        \begin{center}
            \textit{
                Any viable SST extension must include a mechanism—such as an excitation gap—that ensures exponential suppression of internal modes at atomic energy scales.
            }
        \end{center}